\section{进程与调度的实现细节：从 task 到 switch}

前面的章节已经讲了调度算法。本节把进程和调度器按内核实现拆开：进程对象有哪些字段，fork 如何复制资源，exec 如何替换地址空间，wait 如何回收子进程，调度器如何维护 runqueue，睡眠唤醒如何避免丢事件，上下文切换到底切了什么。

\subsection{进程对象之间的引用关系}

Linux 中常用 \code{task\_struct} 表示可调度实体。它不把所有资源都内嵌进去，而是引用地址空间、文件表、信号处理、凭据、命名空间、cgroup、调度实体等对象。这样线程可以共享一部分资源，fork 可以复制或共享不同对象，exec 可以替换地址空间但保留 PID 和部分文件描述符。

\begin{longtable}{|p{0.22\linewidth}|p{0.38\linewidth}|p{0.28\linewidth}|}
\toprule
对象 & 回答的问题 & 生命周期事件 \\
\midrule
\code{task\_struct} & 这个可调度实体是谁、处于什么状态 & fork 创建、exit 释放 \\\tablerowrule
\code{mm\_struct} & 用户地址空间是什么 & fork COW、exec 替换、exit 释放 \\\tablerowrule
\code{files\_struct} & fd 表指向哪些打开文件 & fork 共享或复制、exec close-on-exec \\\tablerowrule
\code{fs\_struct} & cwd、root、umask 等进程文件系统上下文 & chdir/chroot/fork \\\tablerowrule
\code{sighand} & 信号处理函数和 mask & fork/exec/signal \\\tablerowrule
\code{cred} & uid、gid、capability、安全标签 & setuid、exec、LSM 检查 \\\tablerowrule
\code{sched\_entity} & 调度时间与队列状态 & enqueue、dequeue、account runtime \\
\bottomrule
\end{longtable}

分析进程生命周期时要明确引用计数和共享边界：线程共享 \code{mm} 和 \code{files} 时，
一个线程关闭 fd 会影响同进程其他线程；最后一个引用消失时，相应对象才进入释放阶段。

\subsection{fork 的实现顺序}

fork 创建新的 task，并按照资源类型决定复制或共享，随后建立父子关系，把子任务放入 runnable 集合，最后向父子返回不同的值。若中间失败，必须按已经完成的步骤逆序回滚。

\begin{lstlisting}
fork sketch:
  allocate task_struct and kernel stack
  copy scheduler attributes
  copy or share mm_struct with COW PTEs
  copy or share files_struct
  copy signal and credential state
  assign pid
  link into parent children list
  make child runnable
  return child pid to parent, 0 to child
\end{lstlisting}

这里的关键边界是新任务何时对调度器可见。所有必要状态初始化完成前，子任务不能进入 runqueue，否则另一个 CPU 可能立刻调度它并访问尚未就绪的对象。父子关系也要在锁保护下建立，否则 wait 路径可能遗漏子进程。

\subsection{fork 的 COW 页表细节}

fork 复制页表时，匿名可写页通常不复制物理页，而是父子 PTE 都改成只读并标记 COW。物理页引用计数增加。任一方写入时，触发写保护缺页，内核复制新页，再把写入方 PTE 改为可写。

\begin{lstlisting}
fork COW:
  for each writable anonymous PTE:
    clear write bit in parent PTE
    install read-only PTE in child
    mark both as COW
    increment page refcount
  flush parent TLB entries if permissions were tightened

write fault:
  if page refcount == 1:
    make PTE writable
  else:
    allocate new page
    copy old page contents
    decrement old page refcount
    map new page writable
\end{lstlisting}

注意父进程 PTE 也要收紧权限，否则父进程写入不会 fault，就会直接修改共享页，破坏 fork 语义。权限收紧后还要处理 TLB，否则 CPU 可能继续使用旧的 writable 翻译。

\subsection{exec 的关键承诺}

exec 成功后，当前进程身份大体保留，但用户地址空间被新程序替换。PID 不变，部分 fd 保留，设置了 close-on-exec 的 fd 关闭，信号处理按规则重置，凭据可能因 setuid 或文件 capability 改变。

\begin{lstlisting}
execve sketch:
  copy argv/envp from user memory
  open executable and check permissions
  identify binary format
  create new mm_struct
  map program segments
  map interpreter if dynamically linked
  create user stack with argc/argv/envp/auxv
  apply credential transitions
  commit new mm atomically
  start at entry point
\end{lstlisting}

exec 失败要小心。理想情况下，在 commit 新地址空间前失败，旧程序仍能收到错误返回并继续运行；一旦已经不可逆地销毁旧地址空间，就只能终止进程。成熟内核会尽量把可能失败的检查和分配放在 commit 之前。

\subsection{用户栈的构造}

新程序启动时不是 C 的 \code{main} 被内核直接调用。内核把初始用户栈构造成 ABI 规定的布局，包括 \code{argc}、\code{argv} 指针数组、\code{envp} 指针数组、auxiliary vector 和字符串数据。入口点通常是动态链接器或 \code{\_start}。

\begin{lstlisting}
initial user stack high to low:
  strings for argv and envp
  alignment padding
  auxiliary vector entries
  NULL
  envp pointers
  NULL
  argv pointers
  argc
\end{lstlisting}

auxv 给用户态传递页大小、随机数地址、程序头信息、vDSO 入口等。动态链接器用这些信息装载共享库、处理重定位和准备 libc。OS 的 exec 不是“把文件读到内存跳过去”这么简单，它要建立完整 ABI 环境。

\subsection{wait 和 zombie 的意义}

子进程退出后不能立刻完全消失，因为父进程还可能调用 wait 获取退出状态和资源使用信息。于是退出进程进入 zombie 状态，保留 PID、退出码和少量统计；父进程 wait 后，内核释放最后对象。

\begin{lstlisting}
exit:
  close files
  release mm
  record exit_code
  notify parent
  become ZOMBIE
  schedule away forever

wait:
  find child in ZOMBIE state
  copy exit status to parent
  unlink child from parent list
  free final task reference
\end{lstlisting}

若父进程不 wait，zombie 会留下。若父进程先退出，子进程会被 reparent 给 init 或 subreaper。这里的设计保证退出状态不丢，同时避免已退出任务再次运行。

\subsection{runqueue 的最小不变量}

每 CPU runqueue 至少要维护当前任务、runnable 任务集合、锁、调度时钟和统计。一个任务在同一时刻只能处于一个位置：某个 CPU 的 current，某个 runqueue，某个等待队列，或者退出路径。

\begin{rubricbox}
\begin{itemize}
\tightlist
\item RUNNING 任务必须等于某个 CPU 的 current。
\item RUNNABLE 任务必须在且只在一个 runqueue 中。
\item SLEEPING 任务不能留在 runqueue 中。
\item ZOMBIE 任务不能被调度器选中。
\item 迁移任务时，源 runqueue 和目标 runqueue 的锁顺序必须固定。
\item 改变任务状态和队列位置必须在同一临界区内完成。
\end{itemize}
\end{rubricbox}

许多调度错误来自不变量遭到破坏，而非候选选择算法本身：任务重复入队，睡眠任务未出队，唤醒时状态已经变化，或者迁移时两个 CPU 同时观察到同一任务。

\subsection{CFS 的核心数据结构}

公平调度类用 \code{sched\_entity} 记录虚拟运行时间、权重、统计和树节点。runqueue 中的 CFS 队列按 \code{vruntime} 排序，最左边是当前最“亏”的任务，优先运行。

\begin{lstlisting}
CFS account:
  delta_exec = now - curr.exec_start
  weighted = delta_exec * NICE_0_LOAD / curr.weight
  curr.vruntime += weighted
  update rq.min_vruntime

CFS pick:
  next = leftmost entity in rb_tree
  return task_of(next)
\end{lstlisting}

\code{min\_vruntime} 是队列的基准。新唤醒任务不能简单设置为 0，否则会长期占便宜；也不能设置得太大，否则交互响应差。实际实现会把新实体放在相对 \code{min\_vruntime} 合理的位置，并配合唤醒抢占规则。

\subsection{红黑树为什么合适}

公平调度需要频繁插入、删除和取最小 \code{vruntime}。普通链表插入简单但查找最小慢；堆取最小快，但删除任意 sleeping 任务不方便；红黑树能在 \code{O(log n)} 完成插入、删除和按 key 排序。

\begin{longtable}{|p{0.2\linewidth}|p{0.34\linewidth}|p{0.34\linewidth}|}
\toprule
结构 & 优点 & 缺点 \\
\midrule
FIFO 队列 & O(1)，简单 & 不能表达公平权重和 vruntime 排序 \\\tablerowrule
优先级数组 & 固定优先级快 & 动态公平排序不自然 \\\tablerowrule
堆 & 取最小快 & 删除任意实体复杂 \\\tablerowrule
红黑树 & 排序、插入、删除平衡 & 常数成本高于队列 \\
\bottomrule
\end{longtable}

数据结构的选择取决于所需操作。任务会频繁睡眠、唤醒、迁移和改变权重，调度器必须以可接受的成本支持这些变化。

\subsection{睡眠唤醒的无丢失顺序}

等待队列正确性的核心是同一把锁保护“条件检查”和“入队睡眠”。否则唤醒可能发生在检查之后、入队之前，等待者永远睡下去。

\begin{lstlisting}
correct wait:
  lock(obj.lock)
  while !condition:
    prepare_to_wait(waitq, current, SLEEPING)
    unlock(obj.lock)
    schedule()
    lock(obj.lock)
  finish_wait(waitq, current)
  unlock(obj.lock)

waker:
  lock(obj.lock)
  change condition
  wake_up(waitq)
  unlock(obj.lock)
\end{lstlisting}

很多真实内核代码会用更复杂的 helper，但它们都在维护这个原则。条件变量、pipe、socket、futex、completion、wait event 都逃不开同一条规则。

\subsection{上下文切换切的是什么}

上下文切换分成调度器层和架构层。调度器选择 next，更新 current、统计、状态；架构层保存 callee-saved 寄存器、栈指针和必要控制状态，切到 next 的内核栈，再返回到 next 之前停下的位置。

\begin{lstlisting}
schedule:
  prev = current
  next = pick_next_task(rq)
  if prev != next:
    rq.current = next
    context_switch(prev, next)

context_switch:
  switch address space if needed
  switch kernel stack
  save prev callee-saved registers
  restore next callee-saved registers
  return as next task
\end{lstlisting}

\code{context\_switch} 的返回语义需要特别说明：函数返回时，当前执行流已经属于 next 任务。prev 将来再次被调度时，会从它原先离开的位置继续返回。因此，上下文切换在控制流表面仍表现为一次函数返回，返回前后的执行主体却已经改变。

\subsection{地址空间切换和 lazy TLB}

若两个线程共享同一个 \code{mm\_struct}，切换时不需要换页表根。若切到不同进程，通常要切换页表根或地址空间标识。切换页表会影响 TLB；x86-64 PCID 可以减少全量刷新。

\begin{lstlisting}
switch_mm(prev_mm, next_mm):
  if prev_mm == next_mm:
    keep current page table root
  else:
    load next page table root with PCID when enabled
    ensure stale translations cannot be used
\end{lstlisting}

内核线程可能没有自己的用户地址空间，常借用上一个进程的 \code{mm} 或使用特殊内核地址空间。lazy TLB 这类优化的目标，是减少无意义的页表切换和 TLB 刷新。

\subsection{负载均衡和调度域}

多核系统具有层次结构。SMT sibling、同一核心、同一 LLC、同一 NUMA node、跨 socket 的迁移成本不同。调度器通常建立调度域，从近到远进行负载均衡。

\begin{lstlisting}
load balance levels:
  SMT siblings
  cores sharing LLC
  CPUs in same NUMA node
  CPUs across NUMA nodes
\end{lstlisting}

迁移一个任务可能减少 runnable wait，却损失 cache 热状态并增加远端内存访问。x86-64 服务器和工作站上，调度器要估计 SMT sibling、共享 LLC、NUMA node、跨 socket 迁移、turbo 预算和功耗封顶。任务放错位置会让锁竞争、cache miss、远端内存访问和尾延迟一起变差。

\subsection{cgroup CPU 配额}

容器和服务隔离需要限制一组任务的 CPU 使用。cgroup CPU controller 可以按周期给 group 分配 runtime budget。预算耗尽后，该组任务被 throttled，直到下个周期补充预算。

\begin{lstlisting}
cgroup quota:
  period = 100ms
  quota = 20ms
  tasks in group may run total 20ms per period
  after budget exhausted:
    throttle group
  at next period:
    refill runtime and unthrottle
\end{lstlisting}

这会产生一种典型现象：容器内 CPU 使用率尚未达到一个完整核心，请求却出现周期性延迟尖峰，因为 group 已被配额限流。性能分析必须结合 cgroup throttling 统计，宿主机总体 CPU 使用率不足以解释该现象。

\subsection{调度实现测试}

\begin{itemize}
\tightlist
\item fork 后子任务不能在初始化完成前进入 runqueue。
\item COW fork 后父子任一方写入都不能修改另一方可见内容。
\item exec 失败路径不得破坏旧程序继续运行的必要状态。
\item wait 必须能回收 zombie，并且不能回收非子进程。
\item 睡眠和唤醒并发时不能 lost wakeup。
\item 迁移任务时不能同时出现在两个 CPU 的 runqueue。
\item cgroup quota 耗尽时 group 被 throttle，周期补充后能恢复。
\item 上下文切换后 callee-saved 寄存器、栈和地址空间必须正确。
\end{itemize}
